% Part: many-valued-logic
% Chapter: infinite-valued-logics
% Section: goedel

\documentclass[../../../include/open-logic-section]{subfiles}

\begin{document}

\olfileid{mvl}{inf}{god}

\olsection{Logika-logika G\"odel}

\begin{editorial}
  Iki mung draf ringkes saka perangan bab logika G\"odel
  mawa nilai tanpa wates.
\end{editorial}

\begin{defn}\ollabel{def:goedel} Logika G\"odel mawa nilai tanpa wates
~$\LogGod[\infty]$ ditetepake nganggo matriks iki:
\begin{enumerate}
  \item Basa proposisional baku $\Lang L_0$ kang ngemot
  $\lfalse$, $\lnot$, $\land$, $\lor$, $\lif$.
  \item Himpunan nilai kabeneran $V_\infty$.
  \item $1$ mung siji-sijine nilai pinilih, yaiku $V^+ = \{1\}$.
  \item Fungsi kabeneran diwenehake ing ngisor iki:
  \begin{align*}
    \tf{\lfalse} & = 0\\
    \tf{\lnot}[\LogGod](x) & = \begin{cases}
      1 & \text{yen } x =0\\
      0 & \text{yen ora}
    \end{cases}\\
    \tf{\land}[\LogGod](x,y) & = \min(x,y)\\
    \tf{\lor}[\LogGod](x,y) & = \max(x,y)\\
    \tf{\lif}[\LogGod](x,y) & = \begin{cases}
      1 & \text{yen } x \le y\\
      y & \text{yen ora.}
    \end{cases}
    \end{align*}
\end{enumerate}
Logika G\"odel $m$ nilai ditetepake kanthi cara padha, nanging $V = V_m$.
\end{defn}

\begin{prop}
  Logika $\LogGod[3]$ kang ditetepake ing
  \olref[thr][god]{defn:goedel} padha karo $\LogGod[3]$
  kang ditetepake ing \olref{def:goedel}.
\end{prop}

\begin{proof}
  Iki katon kanthi mbandhingake tabel kabeneran konektif ing
  \olref[thr][god]{defn:goedel} karo tabel kang ditemtokake
  dening persamaan ing \olref{def:goedel}:
  \begin{center}
    \begin{tabular}{c|c} 
      $\tf{\lnot}[\LogGod[3]]$ & \\ 
      \hline  
      $1$ & $0$ \\ 
      $1/2$ & $0$ \\
      $0$ & $1$ 
    \end{tabular}
    \quad
    \begin{tabular}{c|ccc} 
      $\tf{\land}[\LogGod]$ & $1$ & $1/2$ & $0$ \\ 
      \hline 
      $1$ & $1$ & $1/2$ & $0$ \\ 
      $1/2$ & $1/2$ & $1/2$ & $0$\\ 
      $0$ & $0$ & $0$ & $0$ 
    \end{tabular}
    \\[2ex]
    \begin{tabular}{c|ccc} 
      $\tf{\lor}[\LogGod]$ & $1$ & $1/2$ & $0$ \\ 
      \hline 
      $1$ & $1$ & $1$ & $1$ \\ 
      $1/2$ & $1$ & $1/2$ & $1/2$ \\
      $0$ & $1$ & $1/2$ & $0$ 
    \end{tabular}
    \quad
    \begin{tabular}{c|ccc} 
      $\tf{\lif}[\LogGod]$ & $1$ & $1/2$ & $0$ \\ 
      \hline 
      $1$ & $1$ & $1/2$ & $0$ \\ 
      $1/2$ & $1$ & $1$ & $0$  \\ 
      $0$ & $1$ & $1$ & $1$ 
    \end{tabular}
  \end{center} 
\end{proof}

\begin{prop}\ollabel{prop:god-infty-m}
  Yen $\Gamma \Entails[\LogGod[\infty]] !B$, mula $\Gamma
  \Entails[\LogGod[m]] !B$ kanggo kabeh~$m \ge 2$.
\end{prop}

\begin{proof}
  Latihan.
\end{proof}

\begin{prob}
  Buktèkna \olref[mvl][inf][god]{prop:god-infty-m}.
\end{prob}

Satemene arah kosok baline uga bener.

Kaya $\LogGod[3]$, $\LogGod[\infty]$ nduweni kabeh
!!{formula}s kang sah ing logika intuisionistik minangka tautologi.
Conto-conto kang padha kang dudu tautologi uga tetep dudu tautologi
ing~$\LogGod[\infty]$:
\begin{align*}
  & p \lor \lnot p && (p \lif q) \lif (\lnot p \lor q) \\
  & \lnot\lnot p \lif p && \lnot(\lnot p \land \lnot q) \lif (p \lor q) \\
  & ((p \lif q) \lif p) \lif p && \lnot(p \lif q) \lif (p \land \lnot q)
\end{align*}
Conto !!{formula} kang ora sah ing logika intuisionistik nanging
tautologi ing~$\LogGod[3]$, yaiku $(p \lif q) \lor (q \lif
p)$, uga tautologi ing~$\LogGod[\infty]$. Malah,
$\LogGod[\infty]$ bisa dicireni minangka logika intuisionistik
kang ditambahi skema $(!A \lif !B) \lor (!B \lif !A)$.
Iki dibuktekake dening Michael Dummett, mula $\LogGod[\infty]$
asring diarani logika G\"odel--Dummett~$\Log{LC}$.

\begin{prob}
  Tuduhna yen $(p \lif q) \lor (q \lif p)$ iku tautologi
  ing~$\LogGod[\infty]$.
\end{prob}

\begin{prob}
  Tuduhna yen $(p \lif q) \lor (q \lif r) \lor (r \lif s)$,
  kang dadi tautologi ing $\LogGod[3]$, dudu tautologi
  ing~$\LogGod[\infty]$.
\end{prob}

\end{document}
